%! Equations and Graphs in Both Directions — Unit 2, Lesson 2.3 — Exit Ticket
\documentclass[10pt]{article}
\usepackage{atda-article}
\usepackage{atda-boxes}
\newcommand{\TallMath}[1]{$\displaystyle #1\rule[-1.4em]{0pt}{3.2em}$}
\setlength{\workrowsep}{8pt}

\begin{document}

\pageheader{Unit 2, Lesson 2.3}{Exit Ticket}
% NAMESTRIP: no \namedateperiod here. The student writes their name once, on the
% cover this component is stapled behind. See references/conventions.md.

\vspace{0.15in}

\begin{notesbox}{Both Directions, and One Substitution}
\small
Notes closed. Three items.

\begin{enumerate}[leftmargin=1.6em, itemsep=7pt, topsep=4pt]

\item \textbf{Equation $\to$ graph.} For $y=(x+2)^2-1$: give the anchor point, complete the table,
then plot your five points on the grid and join them with a smooth curve. The parent is drawn.
\smallskip

Anchor point: \blank{2.6cm}
\smallskip

\renewcommand{\arraystretch}{1.3}
\begin{tabular}{@{}|c|c|c|c|c|c|@{}}
\hline
$x$ & $-4$ & $-3$ & $-2$ & $-1$ & $0$ \\ \hline
$y=(x+2)^2-1$ & \blank{1.0cm} & \blank{1.0cm} & \blank{1.0cm} & \blank{1.0cm} & \blank{1.0cm} \\
\hline
\end{tabular}

\begin{center}
\begin{tikzpicture}
\begin{axis}[
    width=9.2cm, height=4.6cm,
    xlabel={$x$}, ylabel={$y$},
    xmin=-5.4, xmax=3.4, ymin=-2.4, ymax=5.4,
    xtick={-5,-4,-3,-2,-1,0,1,2,3}, ytick={-2,-1,0,1,2,3,4,5},
    grid=both, grid style={linegray!45}, axis lines=left,
    tick label style={font=\tiny}, label style={font=\scriptsize},
    legend entries={{$y=x^2$ (the parent)}},
    legend style={font=\scriptsize, draw=none, fill=none,
      at={(0.5,1.02)}, anchor=south, legend columns=2}]
  \addplot[royal, very thick, domain=-2.2:2.2, samples=80]{x^2};
\end{axis}
\end{tikzpicture}
\end{center}

\item \textbf{Graph $\to$ equation, no picture given.} An exponential function has the flat line
$y=-3$ as its asymptote, and its graph passes through $(0,-2)$. Write its equation, name the two
things you used to get it, and give the range in interval notation.

\writelines{3}

\item A classmate was asked for \emph{the exponential parent moved $3$ units to the right} and typed
\texttt{Y1=2\^{}X-3} into the calculator. Notice that this is the same expression you wrote in item
2. Settle whether it followed the instruction, \textbf{using one substitution}, and give the
keystrokes that would be correct.

\writelines{3}

\end{enumerate}
\end{notesbox}

\end{document}
